% Source: MQ19b_v1.html

\documentclass[11pt]{article}
\usepackage[margin=1in]{geometry}
\usepackage{amsmath}
\usepackage{amssymb}
\usepackage{siunitx}
\usepackage{enumitem}

\title{MQ19b: Heat Transfer}
\author{}
\date{}

\begin{document}
\maketitle

\section*{MQ19b: Heat Transfer}

\begin{enumerate}[leftmargin=*,label=\textbf{Question \arabic*.}]

\item \hfill \textit{1 pts}

A monatomic ideal gas undergoes an isothermal compression from a volume of 2.7 m$^{3}$ and pressure 60 kPa to a final volume of 1 m$^{3}$. Calculate the work in kJ done \textbf{by} the gas on its environment.

\item \hfill \textit{1 pts}

A monatomic ideal gas expands at constant temperature from 1 m$^{3}$ to 2 m$^{3}$. If the initial pressure is 136 kPa, calculate the work in kJ done \textbf{on} the gas during this process.



\emph{W}$_{on}$ = \rule{2.5cm}{0.4pt} kJ

\item \hfill \textit{1 pts}

A monatomic ideal gas is compressed adiabatically starting at a volume of 3 m$^{3}$ and pressure of 196 kPa and ending at a volume of 1.6 m$^{3}$. Find the final pressure in kPa.

\item \hfill \textit{1 pts}

A diatomic ideal gas decreases in pressure adiabatically from an initial pressure of 1.50 × 10$^{5}$ Pa and volume of 1.00 m$^{3}$ to a final volume of 2.00 m$^{3}$.  Two answers are required.

(a) Calculate the final pressure in kPa. Answer:

    [ Select ]

  ["67.2 kPa", "47.4 kPa", "56.8 kPa", "82.3 kPa", "95.6 kPa"]

(b) Calculate the change in the internal energy in kJ. Answer:

    [ Select ]

  ["-47.9 kJ", "-90.8 kJ", "-75.4 kJ", "-58.8 kJ", "-82.5 kJ"]

\item \hfill \textit{1 pts}

A monatomic ideal gas expands adiabatically from 1.00 m$^{3}$ to 2.2 m$^{3}$. If the initial pressure is 247 kPa, calculate the work done \textbf{by} the system in kJ.

\item \hfill \textit{1 pts}

On a \emph{pV}-diagram, a diatomic gas process proceeds along the straight line from (3 m$^{3}$, 260 kPa) to ( 4 m$^{3}$, 487 kPa). Calculate the work done by the gas in kJ.

\end{enumerate}

\end{document}
